% GENERATED, edit the French unit or the translation table
\mxunit{terminale-sources}{Sources and scope}
\textbf{Sources provided.} Barbazo, \emph{Terminale spécialité mathématiques}, Hachette Livre, 2020, ISBN 978-2-01-786619-0; annual progression for Terminale specialty 2026--2027 (15 pages); Terminale specialty syllabus provided (new annex); 2026--2027 almanac of the institution. The textbook supplies the precisely referenced exercises in every chapter except chapters~\ref{ch:13}, \ref{ch:14} and~\ref{ch:12}; the progression supplies the order and the session identifiers.
\textbf{Calendar.} ``Maths spé'' events, from 4 September 2026 to 2 July 2027, read on 5 September 2026. The descriptions are taken as general intentions; the precise content and the prerequisites follow the progression PDF.
\textbf{Distinction between the syllabuses.} The PDF provided presents the new organisation (vocabulary and logic before the domains) and mentions in particular perpendicular planes. It must not be confused with the 2019 annex applicable to 2026--2027. The scope check was redone on the latter; perpendicular planes are flagged as an extension, with no dedicated assessment.\par
\href{https://eduscol.education.gouv.fr/sites/default/files/document/spe246annexe1158907pdf-84159.pdf}{Terminale specialty syllabus: 2019 annex}\par
\href{https://www.education.gouv.fr/bo/2026/Hebdo14/MENE2602919A}{Order of 26 February 2026: application from the start of the 2027--2028 school year}\par
\textbf{Official verification.} Éduscol page ``Programmes et ressources en mathématiques -- voie GT'', consulted on 5 and 6 September 2026: 2019 Terminale syllabus for this year and new syllabus from 2027--2028.\par
\href{https://eduscol.education.gouv.fr/5817/programmes-et-ressources-en-mathematiques-voie-gt}{Éduscol: mathematics syllabuses and resources}
\textbf{Recent examination reference.} Ministerial specialty mathematics paper, Metropolitan France, 17 June 2025: duration 4 h and four compulsory exercises. Used as a guide for write-ups and as reference work: exercise 1, questions 5 and 6; exercise 3. Their statements are to be read in the past paper, their solutions appear in the teacher edition.\par
\href{https://www.education.gouv.fr/sites/default/files/2025-06/baccalaur-at-g-n-ral-2025-math-matiques-440856.pdf}{Ministry: specialty paper of 17 June 2025}
\textbf{National calendar.} The ministry page for the 2027 session was consulted; it does not replace the specific instructions for overseas centres or the individual summons. The booklet therefore sets no local examination date on this basis alone.\par
\href{https://www.education.gouv.fr/reussir-au-lycee/baccalaureat-brevet-cap-le-calendrier-2027-341384}{Ministry: national calendar for the 2027 session}
\textbf{Guides to presentation.} The statements were rewritten by making explicit the objects, domains, hypotheses and expected results, and by separating conjecture from proof. Examples of write-ups consulted:\par
\href{https://www.maths-et-tiques.fr/telech/20SuitesTS2.pdf}{Maths et tiques, Yvan Monka: limits of sequences}\par
\href{https://www.lyceedadultes.fr/sitepedagogique/documents/math/mathTermSpe/ctrle_2020_2021/13_devoir_01_03_2021.pdf}{Lycée d'adultes, Paul Milan: test of 1st March 2021}\par
\textbf{Original material in the booklet.} Notes, activities, exercises, homework, practice papers and assessments are written for this document. The contextualised models are teaching situations, not empirical data. The distribution calendar and the session choices are made explicit so that they can be adapted to the real pace of the class.
\end{document}


